\section{Prompt Templates for Milestone Generation and Stage Judges}
\label{app:prompts}

This section documents the primary prompt templates used across DynSTEER's pipeline, formatted in structured boxed environments.

\begin{tcolorbox}[colback=gray!4,colframe=black!75,arc=4pt,boxrule=0.8pt,left=8pt,right=8pt,top=6pt,bottom=6pt]
\textbf{Prompt Template 1: Milestone Graph Candidate Generation Prompt} \\[-6pt]
\rule{\textwidth}{0.4pt}
\small
\textbf{System Prompt:} \\
You generate candidate milestone goal graphs, not step-by-step trajectories. A node belongs in one candidate graph only when, under that candidate's complete and feasible interpretation, every valid completion must reach that operation or goal. One candidate is a hypothesis about necessity; necessity across candidates is decided later by code aggregation.

\textbf{User Prompt Input:} \\
This is candidate batch \texttt{\{batch\_index\}}. Focus code: \texttt{\{focus\_code\}}. \\
Batch focus: \texttt{\{focus\_instruction\}}

\textbf{Core Requirements:}
\begin{itemize}[leftmargin=*,itemsep=1pt,topsep=1pt]
  \item Return standalone candidate graphs for the current task based strictly on the public task view \texttt{\{task\_view\}}.
  \item Each graph must define its own milestones, dependency edges, dynamic bindings, and fatal minefields.
  \item Seek genuine diversity: explore distinct valid tools, direct public state lookups, or alternative precedence orders.
  \item Output must strictly follow the JSON schema specifying nodes, edges, dependencies, and minefield operations.
\end{itemize}
\end{tcolorbox}

\begin{tcolorbox}[colback=gray!4,colframe=black!75,arc=4pt,boxrule=0.8pt,left=8pt,right=8pt,top=6pt,bottom=6pt]
\textbf{Prompt Template 2: Standard Judge Unified Multi-Field Prompt} \\[-6pt]
\rule{\textwidth}{0.4pt}
\small
\textbf{System Role:} \\
You are a strict task-trajectory evaluator. Evaluate only the supplied task, stage interval, constraint checks, and trajectory steps. Do not invent external facts.

\textbf{Context Inputs:}
\begin{itemize}[leftmargin=*,itemsep=1pt,topsep=1pt]
  \item \texttt{stage\_goal}: Success condition for the current dominator-bounded stage interval.
  \item \texttt{rubrics}: Target dimensions to score simultaneously (e.g., correctness, efficiency, safety, tool compliance).
  \item \texttt{steps}: Observed actions and tool returns strictly within the stage evidence window $\mathcal{W}(m^*)$.
\end{itemize}

\textbf{Evaluation Instructions:}
\begin{itemize}[leftmargin=*,itemsep=1pt,topsep=1pt]
  \item Determine whether this stage satisfies \texttt{stage\_goal} within the supplied interval.
  \item Score each target dimension from 0.0 to 1.0 based on explicit behavioral evidence.
  \item Provide concise diagnostic justifications citing specific step indices and tool call names.
  \item Output format: Strict JSON mapping dimension names to score, confidence, and diagnostic rationale.
\end{itemize}
\end{tcolorbox}

\begin{tcolorbox}[colback=gray!4,colframe=black!75,arc=4pt,boxrule=0.8pt,left=8pt,right=8pt,top=6pt,bottom=6pt]
\textbf{Prompt Template 3: Expensive Judge Specialist Prompt (Tool Quality Dimension)} \\[-6pt]
\rule{\textwidth}{0.4pt}
\small
\textbf{Specialist Objective:} \\
You are the specialist LLM-Judge for the \texttt{tool\_quality} dimension. Score only tool-use quality.

\textbf{Detailed Forensic Focus:}
\begin{itemize}[leftmargin=*,itemsep=1pt,topsep=1pt]
  \item Whether the agent selected the optimal tools for the active stage goal.
  \item Whether tool arguments are complete, accurate, and type-conforming without hallucinations.
  \item Whether return values, exceptions, and empty results were correctly parsed and handled rather than ignored.
  \item Penalize redundant retry loops, irrelevant exploratory calls, and unhandled tool exceptions.
  \item Output: Calibrated continuous score in $[0, 1]$, confidence score, and forensic audit trace.
\end{itemize}
\end{tcolorbox}
